\section{Introduction}
\label{sec:intro}

Industrial Processes and Product Use (IPPU) is not one calculation. It includes carbonate
decomposition in cement, lime and glass; feedstock carbon in ammonia and petrochemicals;
nitrous oxide from nitric acid; carbon inputs to metals; non-energy uses of fuels; and
fluorinated gases held in equipment and products \citep{ipcc2006ippu,ipcc2019}. These sources
share an inventory sector but not an activity unit, emission factor or spatial representation.
A credible national account must therefore do two things at once: connect each process to its
own physical activity and identify explicitly which categories remain unestimated because
process-specific activity data are insufficient.

Earlier work independently estimated cement calcination from one cement statistic, one
assumed clinker fraction and one emission factor, but did not yet provide a sector-wide IPPU
framework. It also omitted a required term in the cement method. When clinker production is inferred from cement,
IPCC's Tier~1 worksheet requires clinker imports and exports because calcination emissions
occur where clinker is produced rather than where it is ground or consumed. Pakistan is a
substantial net clinker exporter, so that correction is material.

This study asks one question: \emph{how far can a reproducible, physically coherent IPPU
inventory be built from currently public Pakistani evidence while restricting emission
estimates to categories for which the relevant process-specific activity data are available?}

We address three questions. First, what does a correct period- and trade-matched public-data
calculation imply for cement calcination? Second, how closely does inferred national clinker
agree with company-reported production, and which evidence still prevents a complete Tier~2
balance? Third, what does Pakistan's official reporting imply about the material categories
that remain outside the independent estimate? We answer them by reconciling cement production
and trade in a material balance, comparing inferred clinker with producer filings, auditing
the CRT at cell level and assessing data completeness across 2.A--2.G.

The analysis establishes both an estimate and the present limits of the available evidence.
Cement 2.A.1 is estimated for FY2023--24 and reconstructed independently for calendar 2022.
Public data identify glass, soda ash, steel products, lubricants, HFC consumption and partial
ammonia/nitric-acid activity, but estimating their emissions requires additional evidence on
process routes, recovery, abatement or equipment stocks.
